\documentclass[10pt,twocolumn,letterpaper]{article}
\usepackage{cvpr}
\usepackage{times}
\usepackage{graphicx}
\usepackage{amsmath,amssymb,booktabs}
\usepackage[pagebackref,breaklinks,colorlinks]{hyperref}
\input{generated/study.tex}
\title{Noise-Aware Edge Regularization for Few-Step Diffusion Distillation}
\author{Le Duc Dong \quad Nguyen Huu Thuat \quad Pham Viet Hoang \quad La Nguyen Gia Hy\\CE7454 Project 2}
\begin{document}
\maketitle

\begin{abstract}
We study the quality--cost trade-off of diffusion distillation using
\StudyDataset{}. Our proposed loss weights image-gradient supervision by
signal level. Controlled comparisons include teacher shortcuts, progressive
distillation, constant-edge supervision, SwiftBrush-inspired VSD, consistency
distillation, DMD-inspired training, and InstaFlow-inspired reflow. This working draft imports measured
tables and figures; empirical conclusions remain pending.
\end{abstract}

\section{Problem and research question}
Can a few-step student retain more of a teacher's image quality at an identical
training budget by concentrating edge supervision on lower-noise intervals?
\ifStudySD
Our setting is text-to-image generation at $512\times512$ with Stable Diffusion
1.5~\cite{rombach2022ldm}. Rank-\StudyLoRARank{} LoRA adapters~\cite{hu2022lora}
reduce training and checkpoint cost. The pretrained backbone, text encoder,
and variational autoencoder (VAE) remain frozen.
\else
Our setting is unconditional generation of $\StudyResolution\times\StudyResolution$ RGB
\StudyDataset{} images.
\fi
We measure quality, denoiser evaluations (NFE), and synchronized latency
separately to account for hardware and batch effects.

Figure~\ref{fig:samples} presents fixed-noise samples at the indicated inference budgets.
\input{generated/qualitative_figure.tex}

\section{Related work}
DDPM learns denoising predictions at multiple noise levels~\cite{ho2020ddpm}.
DDIM permits deterministic trajectories through selected levels~\cite{song2021ddim}.
Progressive distillation trains one student transition to reproduce two teacher
transitions, then repeats the step-count reduction~\cite{salimans2022progressive}.
Consistency models pursue an alternative endpoint-prediction
objective~\cite{song2023consistency}. Distribution matching distillation instead
matches generated and teacher distributions using score
information~\cite{yin2024dmd}. SwiftBrush uses variational score distillation
without external training images~\cite{nguyen2024swiftbrush}.
InstaFlow straightens noise--image couplings with reflow before one-step
distillation~\cite{liu2024instaflow}. Recent work includes WaDi's low-rank
weight-direction adapters~\cite{wang2026wadi} and pixel MeanFlow's one-step
generation without a latent autoencoder~\cite{lu2026pixelmeanflow}. These
2026 methods provide context rather than measured baselines: their adapter or
pixel-space model settings differ from our fixed SD1.5/LoRA protocol.
We evaluate the regularizer as a hypothesis without claiming established novelty.

\section{Method}
\paragraph{Forward process and deterministic transitions.}
Let $x_t=\alpha_t x_0+\sigma_t\epsilon$, where
$\alpha_t^2=\bar\alpha_t$, $\sigma_t^2=1-\bar\alpha_t$, and
$\epsilon\sim\mathcal N(0,I)$. The network predicts noise, with clean-image
estimate $\hat x_0=(x_t-\sigma_t\epsilon_\theta(x_t,t))/\alpha_t$.
An unclipped deterministic DDIM step to $s<t$ is
\begin{equation}
 x_s=\alpha_s\hat x_0+\sigma_s\epsilon_\theta(x_t,t).
\end{equation}
The clean boundary uses $\alpha=1,\sigma=0$. Student grids take every second
teacher point, preserving endpoints between training and inference.
\ifStudySD
Here $x$ denotes $4\times64\times64$ scaled VAE latents, and denoisers receive
frozen CLIP text embeddings. The first teacher uses classifier-free guidance
(CFG) \StudyGuidance{}; later teachers and students use one conditional branch.
Guidance is incorporated into distillation targets rather than applied again to
the student. A guided teacher step counts as two NFE even when batched together.
\fi

\paragraph{Two-to-one distillation.}
The frozen teacher takes two steps, $t\rightarrow m\rightarrow s$, to produce
$x_s^T$. Inverting one student step yields a stopped clean-image target
\begin{equation}
 \tilde x_0=\frac{x_s^T-(\sigma_s/\sigma_t)x_t}
 {\alpha_s-(\sigma_s/\sigma_t)\alpha_t}.
\end{equation}
Each stage's EMA initializes the next teacher and student. Our baseline adapts
progressive distillation with pretrained epsilon prediction and clipped-SNR
weighting, rather than the original parameterization and schedule.

\paragraph{Proposed edge loss.}
Write $D_x,D_y$ for first-order horizontal and vertical image differences.
Per image, let $E_\nabla$ be the average squared discrepancy between these
differences in $\hat x_0$ and $\tilde x_0$. We optimize
\begin{equation}
 \mathcal L=\mathbb E\!\left[w_t\left(
 \|\hat x_0-\tilde x_0\|^2+\lambda g_t E_\nabla\right)\right],
\end{equation}
where $w_t=\min(\mathrm{SNR}(t),5)$ and $g_t=\alpha_t^2$.
Reconstruction and edge terms use mean reductions. \ifStudySD
The reconstruction term uses latent space; $E_\nabla$ uses decoded RGB
predictions and targets. The frozen decoder passes gradients to the student,
with activation checkpointing to limit memory.\fi{}
The baseline sets $\lambda=0$;
the constant-edge ablation sets $g_t=1$; the proposed variant uses
$\lambda=0.1$ and the noise-aware gate, targeting reliable spatial signal.
Over-sharpening and reduced diversity remain possible.

\paragraph{SwiftBrush-inspired comparison.}
\ifStudySD
We adapt score distillation to caption-conditioned SD latents.
\else
We adapt score distillation to unconditional \StudyDataset{} pixels.
\fi
Let $G_\theta(z)$ be the one-step student's clean estimate. A frozen teacher
$\epsilon_\phi$ supplies the real score and a trainable auxiliary U-Net
$\epsilon_\psi$ estimates the current generator's score. For
$y_t=\alpha_tG_\theta(z)+\sigma_t\epsilon$, we inject the stopped gradient
$\sigma_t^2(\epsilon_\phi(y_t,t)-\epsilon_\psi(y_t,t))$ into the generator.
The auxiliary model minimizes $\alpha_t^2\|\epsilon_\psi(y_t,t)-\epsilon\|^2$
on detached images, alternating one generator and one scorer update.
\ifStudySD
Teacher and auxiliary predictions use the same CFG \StudyGuidance{}; the
auxiliary denoising update drops text conditioning with probability 0.1.
Both networks use rank-\StudyLoRARank{} adapters, replacing SwiftBrush's full
generator training. Our optimizer, learning rates, EMA and scorer sample reuse
also differ from the source recipe; this is \emph{SwiftBrush-inspired VSD}.
A separate settings variant retains rank-32 students with source Adam rates,
rank-64 q/v scorers and CFG 4.5; its results are excluded from this shared-CFG table.
\else
We replace latent text conditioning, CFG, and LoRA with unconditional pixels
and a full auxiliary U-Net. This is labelled \emph{SwiftBrush-inspired VSD}.
\fi

\ifStudySD
\paragraph{Consistency distillation.}
We match a student's endpoint prediction at $t$ to a stopped EMA endpoint at
$s$, after one guided teacher DDIM transition. For $r=\sigma_t/\alpha_t$ and
$b=0.5$, the endpoint is
\begin{equation}
 f_\theta(x_t,t)=\frac{b^2x_t/\alpha_t}{r^2+b^2}
                  +\frac{r\hat x_0}{\sqrt{r^2+b^2}}.
\end{equation}
This enforces $f(x,-1)=x$. Adjacent intervals use the original 64-point grid;
latent L2 trains one-step inference. This VP/LoRA blend adapts consistency
distillation~\cite{song2023consistency}, rather than its EDM/LPIPS recipe.

\paragraph{DMD-inspired comparison.}
We normalize $\hat x_0^{F}-\hat x_0^{T}$ by the per-sample mean absolute residual
$|G-\hat x_0^{T}|$ and inject the stopped direction into $G$~\cite{yin2024dmd}.
The teacher uses CFG \StudyGuidance{}; the fake scorer uses CFG 1 and unweighted
epsilon denoising. A separate paired stream adds $0.25$ times latent L2 to fixed
noise--teacher-output pairs (\StudyDMDPairs{} per seed). These use 64 guided
DDIM steps. LoRA and latent L2 replace full-generator training and perceptual
regression. This is DMD-inspired; pair time and NFE are reported.

\paragraph{InstaFlow-inspired comparison.}
We cache the same type of noise--teacher-output pairs, then fit a direct
velocity head on $x_u=(1-u)z+u x_T$, with target $x_T-z$ and uniform $u$.
The first half of the generator budget trains reflow. Its frozen EMA generates
a new coupling using 64 Euler steps; the second half fits its endpoint velocity
at $u=0$. Inference is $x=z+v_\theta(z,0,c)$: one step, one NFE.
\ifStudyFlowVP
The velocity is $v_\theta=\hat x_0-\epsilon_\theta$, with normalized input
$x_u/\sqrt{u^2+(1-u)^2}$ and the nearest pretrained VP signal-to-noise level.
This conversion preserves the frozen epsilon head's meaning; the nonzero
terminal signal remains an approximation at $u=0$.
\else
The U-Net predicts velocity directly with time labels
$\operatorname{round}(999(1-u))$.
\fi
Rank-matched LoRA and \StudyFlowPairs{} cached pairs per phase adapt
InstaFlow~\cite{liu2024instaflow}; this is not its released
full-model training recipe. Both pair-generation passes count toward training
time and the reported pair NFE.
\fi

\section{Experimental protocol}
The default teacher is the public \texttt{\StudyTeacher} checkpoint
at the pinned revision recorded in the configuration.
\ifStudySD
COCO-2017~\cite{chen2015coco} images are sorted by image ID;
\ifStudyContinuation we skip the original \StudyTrainOffset{} training images and
use the remaining \StudyTrainCount{} fresh images with shuffled full epochs.
\else the first \StudyTrainCount{} training images are selected.\fi
Each image uses its lowest-ID caption. Images use
bicubic resize and center crop to 512 pixels, with no flips that could contradict
captions. Evaluation uses the first \StudyReferenceCount{} validation image--caption
pairs under the same ordering. These are fixed subsets rather than full-dataset
benchmark results.
\else
Progressive training uses only the
\StudyTrainCount{}-image training split, with random horizontal flips and uniform
minibatch sampling with replacement. Final quality evaluation uses the official
\StudyReferenceCount{}-image \StudyReferenceSplit{} split.
\fi
Pilot selection must be disclosed before the final sweep. \StudySplitNote{}

\ifStudyContinuation
Each method resumes its completed one-step student, optimizer, and EMA state
after \StudyParentUpdates{} generator updates. The continuation uses only
one-step generation and adds \StudyStageUpdates{} generator updates per method,
bringing the cumulative budget to \StudyTotalUpdates{}. Effective batch size is
\StudyBatchSize{} (\StudyMicroBatch{} per microbatch, accumulating
\StudyAccumulation{} microbatches); optimizer families and loss settings are
preserved, with a new cosine schedule for AdamW methods. Progressive variants
retain their frozen two-step teachers. Consistency retains its adjacent teacher
grid; DMD preserves its scorer and uses fresh regression pairs. InstaFlow
continues only endpoint distillation with newly coupled prompts from the
parent's frozen reflow teacher; its reflow optimization phase is not repeated.
SwiftBrush's CFG-4.5 source-recipe adaptation is evaluated separately from this
CFG-\StudyGuidance{} cohort. Training seed(s) are \StudySeeds{}; evaluation
captions and initial noise are shared. Auxiliary updates and pair-generation
cost are reported separately from generator updates.
\else
The configured curriculum starts with \StudyInitialSteps{} teacher steps:
$\StudyCurriculum$.
Progressive variants receive \StudyStageUpdates{} generator updates per stage,
effective batch size \StudyBatchSize{} (\StudyMicroBatch{} per microbatch, accumulating
\StudyAccumulation{} microbatches), AdamW with learning rate $\StudyLearningRate$,
and EMA decay \StudyEMADecay{}.
Training seed(s) are \StudySeeds{}. Every method is evaluated using
the same initial evaluation noise\ifStudySD{} and captions\fi. The VSD adaptation receives \StudyVSDUpdates{} generator
updates, matching the cumulative progressive budget at one step, plus \StudyVSDUpdates{}
auxiliary updates. It consumes no real images during distillation\ifStudySD{}, but uses the same training captions\fi. Generator
update matching does not equalize compute; scorer updates and time
are reported separately. \ifStudySD
Consistency and DMD also receive \StudyVSDUpdates{} generator updates;
DMD adds \StudyVSDUpdates{} scorer updates and a separate regression stream.
InstaFlow splits the same total generator updates equally between reflow and
endpoint distillation and adds two disclosed pair-generation passes.
EMA tracks adapter weights.\else The optional scratch teacher is separate.\fi
\fi
\ifStudySD The matched one-step pilot trains every student through one-step inference;
teacher rows are references with separately stated NFE. Optional edge-weight
ablations use $\lambda\in\{0.03,0.1,0.3\}$ with the same seeds and updates,
alongside no-edge and constant-edge controls.\fi

\paragraph{Imported evaluation settings.}
\input{generated/protocol.tex}
FID and Inception Score use the same pinned \texttt{torch-fidelity}
implementation. \ifStudySD CLIP ViT-B/32 cosine similarity measures text alignment;
we report $100\max(\cos(I,c),0)$ averaged over image--caption pairs.\fi{}
Cross-paper FID rankings are not supported by this protocol.
Latency excludes disk I/O, PNG export, metric extraction, and initial noise
allocation, includes all denoiser transitions\ifStudySD{}, text encoding, and VAE decoding\fi, and uses warm-up plus device
synchronization. Runtime rows are separated by hardware and batch size.

\section{Results and analysis}
\paragraph{Current evidence.}
\input{generated/result_status.tex}
\IfFileExists{generated/checkpoint_evidence.tex}{\input{generated/checkpoint_evidence.tex}}{}
\begin{table}[t]
 \centering\scriptsize
 \resizebox{\linewidth}{!}{\input{generated/quality_table.tex}}
 \caption{One-step student quality and teacher references under the shared protocol. Values with $\pm$ denote
 variation across training seeds. Single-seed rows do not establish robustness.}
 \label{tab:quality}
\end{table}
\begin{table}[t]
 \centering\scriptsize
 \resizebox{\linewidth}{!}{\input{generated/latency_table.tex}}
 \caption{Synchronized latency, grouped by actual hardware and batch size.
 Compare speed only within the same hardware and precision setting.}
 \label{tab:latency}
\end{table}
\begin{table}[t]
 \centering\scriptsize
 \resizebox{\linewidth}{!}{\input{generated/cost_table.tex}}
 \caption{Cumulative one-step training cost. G/F count generator/scorer updates;
 Pair NFE counts guided-teacher and reflow-teacher calls to build cached pairs. Hours include
 pair generation. Matching G does not match total compute.}
 \label{tab:cost}
\end{table}

\input{generated/ablation_section.tex}
\input{generated/extra_qualitative.tex}

\paragraph{Analysis to complete after training.}
Assess the gate against both controls across seeds and inference budgets.
Inspect prompt/noise-matched failures, diversity, stability, memory, and time;
account for scorer and regression costs. Conclusions require measured evidence.

\section{Limitations and conclusion}
One teacher and a limited dataset restrict generality.
\ifStudySD
Small adapters and training budgets may restrict quality.
\else
It cannot establish performance on high-resolution text-to-image generation.
\fi
Epsilon parameterization and nonzero terminal signal may limit one-step quality.
FID depends on sample count and incompletely describes diversity and artifacts.
The gate is weakest at high noise; one-step negative results are informative.
This unfinished draft claims no empirical improvement.

{\small
\bibliographystyle{ieeenat_fullname}
\bibliography{references}
}
\end{document}
